\documentclass[a4paper]{article}

% Basic packages
% Español (México) con XeLaTeX: fontspec para fuentes Unicode, polyglossia
% para la división silábica y los nombres del documento en la variante mexicana.
\usepackage{fontspec}
\usepackage{polyglossia}
\setdefaultlanguage[variant=mexican]{spanish}
% Los nombres chinos van como «pinyin (original)»: xeCJK da a los caracteres
% Han su propia fuente; sin ella XeLaTeX los omite con un aviso.
% FandolSong no cubre algunos caracteres de nombres (U+73FA); WenQuanYi Zen Hei sí.
\usepackage[AutoFallBack=true]{xeCJK}
\setCJKmainfont{FandolSong-Regular.otf}
\setCJKfallbackfamilyfont{\CJKrmdefault}{WenQuanYi Zen Hei}
% polyglossia no traduce los nombres de listings.
\AtBeginDocument{\providecommand{\lstlistingname}{}\renewcommand{\lstlistingname}{Listado}%
  \providecommand{\lstlistlistingname}{}\renewcommand{\lstlistlistingname}{Índice de listados}}
\usepackage{amsmath, amssymb}
\usepackage{graphicx}
\usepackage[margin=2.5cm]{geometry}
\usepackage[most]{tcolorbox}
\usepackage{etoolbox}
\usepackage{listings}
\usepackage{hyperref}
\usepackage{booktabs}
\usepackage{subcaption}
\usepackage{float}

% Highlight boxes
\newtcolorbox{knowledgebox}[1]{
    enhanced, colback=blue!5!white, colframe=blue!75!black, colbacktitle=blue!75!black,
    coltitle=white, fonttitle=\bfseries, title={#1}, attach boxed title to top left={yshift=-2mm, xshift=2mm},
    boxrule=1pt, sharp corners
}

\newtcolorbox{importantbox}[1]{
    enhanced, colback=yellow!10!white, colframe=yellow!80!black, colbacktitle=yellow!80!black,
    coltitle=black, fonttitle=\bfseries, title={#1}, sharp corners
}

\newtcolorbox{warningbox}[1]{
    enhanced, colback=red!5!white, colframe=red!75!black, colbacktitle=red!75!black,
    coltitle=white, fonttitle=\bfseries, title={#1}, sharp corners
}

% Listings
\lstset{
    language=Python,
    basicstyle=\ttfamily\small,
    keywordstyle=\color{blue},
    stringstyle=\color{red!60!black},
    commentstyle=\color{green!60!black},
    breaklines=true,
    frame=single,
    numbers=left,
    numberstyle=\tiny\color{gray},
    captionpos=b,
    extendedchars=true
}

% Front-page metadata
\newcommand{\notetitle}{Weng Jiayi (翁家翌): OpenAI, GPT, aprendizaje por refuerzo, Infra, posentrenamiento, Tianshou (天授), tuixue, código abierto, CMU, Tsinghua}
\newcommand{\noteauthors}{Elaboradas a partir de materiales públicos del curso}
\newcommand{\notedate}{2026-01-17}
\newcommand{\videochannel}{WhynotTV}
\newcommand{\videopublishdate}{2026-01-17}
\newcommand{\videoduration}{122 minutos}
\newcommand{\videourl}{https://www.bilibili.com/video/BV1darmBcE4A/}
\newcommand{\videocoverpath}{cover.jpg}
\newcommand{\repourl}{https://github.com/hqhq1025/ai-course-notes}

% Footer with repo info
\usepackage{fancyhdr}
\pagestyle{fancy}
\fancyhf{}
\fancyfoot[C]{\thepage}
\fancyfoot[R]{\footnotesize\color{gray}\href{\repourl}{AI Course Notes}}
\renewcommand{\headrulewidth}{0pt}
\renewcommand{\footrulewidth}{0pt}

\begin{document}

\begin{titlepage}
\centering
{\Large Notas de la entrevista\par}
\vspace{1.2cm}
{\huge\bfseries \notetitle\par}
\vspace{0.8cm}
{\large \noteauthors\par}
\vspace{0.3cm}
{\large \notedate\par}
\vspace{1.2cm}

\ifdefempty{\videocoverpath}{
{\small Al generar las notas, indique la ruta local de la portada del video.\par}
}{
\includegraphics[width=0.82\textwidth,height=0.45\textheight,keepaspectratio]{\videocoverpath}\par
}

\vfill
\begin{tcolorbox}[width=0.9\textwidth, colback=black!2!white, colframe=black!60, sharp corners]
\textbf{Autor o canal del video}: \videochannel\par
\textbf{Fecha de publicación}: \videopublishdate\par
\textbf{Duración del video}: \videoduration\par
\ifdefempty{\videourl}{}{
\textbf{Enlace del video}: \href{\videourl}{\nolinkurl{\videourl}}\par
}\par
\textbf{Página del proyecto}: \href{\repourl}{\nolinkurl{\repourl}}\par
\end{tcolorbox}
\end{titlepage}

\tableofcontents
\newpage

%% --- Inicio del contenido --- %%

\section{Presentación del invitado}

Jiayi Weng (翁家翌) ingresó en 2016 a la licenciatura en Ciencias de la Computación de la Universidad de Tsinghua, en 2020 partió a la Universidad Carnegie Mellon (CMU) a cursar la maestría y en 2022 se incorporó a OpenAI. Es uno de los constructores centrales de \textbf{Post-Training RL Infrastructure} dentro de OpenAI: detrás de cada modelo grande que OpenAI ha publicado, desde ChatGPT (GPT-3.5) y GPT-4o hasta GPT-5, está su nombre.

Antes de unirse a OpenAI, Jiayi Weng ya había generado un impacto amplio a través de proyectos de código abierto: liberó en Tsinghua todas sus tareas y materiales de curso para romper la asimetría de información; creó el framework de aprendizaje por refuerzo \textbf{Tianshou (天授)}, que se convirtió en uno de los frameworks ligeros más populares de la comunidad de RL; y durante la pandemia desarrolló el sistema gratuito de consulta de visas \textbf{tuixue.online}, que acumuló millones de visitas.

Esta edición del podcast WhynotTV es una conversación a fondo de más de dos horas que arranca con la infancia de Jiayi Weng y abarca sus estudios, el código abierto, sus decisiones profesionales, la forma de trabajar dentro de OpenAI, y sus reflexiones filosóficas sobre el futuro de la IA y el sentido de la vida.

\begin{knowledgebox}{Contexto de la entrevista}
El outline de esta edición del podcast lo preparó el conductor con la función Deep Research de GPT-5, y Jiayi Weng es precisamente uno de los desarrolladores centrales detrás de GPT-5. El conductor lo llamó «un ciclo maravilloso».
\end{knowledgebox}

\subsection{Resumen de la sección}

Las tres palabras clave de Jiayi Weng son: \textbf{aprendizaje por refuerzo, Post-Training, Infra}. Pero su historia va mucho más allá de lo técnico: es alguien que hace «caridad» con código, que cree en el código abierto y en la igualdad de acceso a la información, y que mantiene un pensamiento independiente en el centro de la tormenta mundial de la IA.

\section{Crecimiento y educación: de las Olimpiadas de Matemáticas a la OI}

\subsection{Talento matemático y aprendizaje autodirigido}

Weng Jiayi (翁家翌) empezó a estudiar para las olimpiadas de matemáticas desde primer año de primaria y pronto mostró talento en esa área. Describe así su característica: \textbf{aprende cosas nuevas despacio, pero una vez que las entiende las usa con extrema rapidez}. Explica este fenómeno con la metáfora de un árbol de conocimiento: otras personas necesitan partir de la raíz y avanzar rama por rama, nivel por nivel, hasta llegar a la conclusión, mientras que él construye directamente un shortcut (atajo) que se salta la deducción intermedia y llega directo a la respuesta.

\begin{importantbox}{Metodología de aprendizaje: aprendizaje lento + aplicación rápida}
Weng Jiayi cuenta que necesita entre dos y tres veces más tiempo que otras personas para aprender algo nuevo, y que también lee código más despacio. Pero su estrategia para compensarlo es: \textbf{aprender por adelantado}. Terminó las matemáticas de preparatoria en segundo de secundaria y empezó cálculo en tercero de secundaria. A esta conducta la llama «invertir en el futuro», una idea que atraviesa toda su carrera profesional.
\end{importantbox}

Esta manera de pensar en «invertir en el futuro» no vino de una exigencia de sus padres, sino que fue completamente espontánea. Su lógica central es: en lugar de desperdiciar el tiempo en resolver ejercicios del momento, más vale aprender algo útil para el futuro, porque el beneficio posterior será mayor. Vale la pena notar que esta manera de pensar a largo plazo ya se había formado desde su época de secundaria.

Weng Jiayi también compartió un detalle sobre su método de estudio: de niño, cuando memorizaba textos escolares, antes de dormir intentaba todos los métodos posibles y recitaba el texto completo de manera entrecortada, con muchas pausas, pero insistía en terminarlo. Después dormía, y al despertar al día siguiente descubría que podía recitarlo de corrido sin ningún esfuerzo. Compara este proceso con la relación entre el Sistema 1 y el Sistema 2: dicho en términos de moda, un cálculo mental pasa directo por el Sistema 1, mientras que entender un conocimiento nuevo requiere que el Sistema 2 lo construya poco a poco, pero una vez construido se convierte en un reflejo del Sistema 1.

\subsection{Iniciación en la programación y la competencia de OI}

Weng Jiayi tuvo su primer contacto con la programación desde primero de secundaria, gracias a un club de interés de programación en su escuela. En preparatoria, por la presión de admisión universitaria, empezó a participar en la Olimpiada de Informática (OI). Reconoce que, al no ser originario de Beijing, sin competir era «tan difícil como subir al cielo» entrar a Tsinghua.

En su camino en la OI, atravesó dificultades:
\begin{itemize}
    \item En primero de preparatoria, en la selección provincial casi no pudo resolver ningún problema.
    \item Gracias a un problema de cobertura mínima con dos elementos, obtuvo la puntuación más alta de todo el concurso y logró, con dificultad, entrar al equipo provincial de Fujian.
    \item En el concurso nacional (NOI), quedó \textbf{en último lugar} dentro del equipo provincial de Fujian, y obtuvo una medalla de bronce.
    \item Renunció a la admisión directa, más segura, en la línea de primer nivel de la Universidad Jiao Tong de Shanghái, y eligió la admisión condicionada de Tsinghua con una rebaja de 60 puntos.
\end{itemize}

\begin{knowledgebox}{Ruta de admisión por competencias de OI (en esa época)}
La ruta típica de admisión por competencias era: NOIP (concurso provincial del grupo de mejora) $\rightarrow$ selección provincial $\rightarrow$ campamento de verano de Tsinghua/Peking (que podía otorgar una rebaja de puntos o admisión directa) $\rightarrow$ concurso nacional NOI. La medalla de oro daba admisión directa; a partir de la medalla de plata se obtenía la condición de admisión en la línea de primer nivel; una rebaja de 60 puntos requería que el resultado del examen de admisión a la universidad (gaokao) estuviera a la altura.
\end{knowledgebox}

Durante tercero de preparatoria, Weng Jiayi siguió estudiando OI en secreto, e incluso llegó a dominar la habilidad de escribir y enviar código directamente en el Safari de un iPad, sin compilador ni editor de código. Considera que esta experiencia entrenó enormemente su capacidad de tener una comprensión completa de la lógica de todo un programa y su capacidad de reacción para localizar errores con rapidez.

También le apasionaba la \textbf{optimización de constantes}: con la misma complejidad algorítmica en tiempo, optimizar el coeficiente constante de la ejecución del código y la longitud del código. El sistema de evaluación de la OI ordenaba los casos de prueba según el tiempo de ejecución, y quien corría más rápido quedaba en primer lugar; si el tiempo de ejecución era el mismo, se ordenaba por longitud de código. Weng Jiayi optimizaba (optimize) ambos indicadores a la vez: «aunque no sirve de mucho, es muy interesante». Esta búsqueda de una optimización extrema del código sentó, más adelante, las bases para construir la infraestructura de RL (RL Infra) en OpenAI.

\subsection{La decisión del gaokao: apostar en medio de la incertidumbre}

Después de obtener la medalla de bronce en el NOI, Weng Jiayi enfrentó una decisión difícil: la condición que ofrecía el campamento de verano de Tsinghua era una \textbf{rebaja de 60 puntos} (sumar 60 puntos al resultado del gaokao; con pasar la línea de corte quedaba admitido), pero con una condición adicional: si en el NOI lograba entrar entre los primeros 150 lugares (línea de medalla de plata), obtendría la admisión directa en la línea de primer nivel. Al final no consiguió la medalla de plata.

La otra opción, más segura, era la admisión en la línea de primer nivel de la Universidad Jiao Tong de Shanghái. Weng Jiayi recuerda que, en ese entonces, durante el segundo semestre de segundo de preparatoria no había estudiado las materias generales y no sabía cuántos puntos podría obtener en el gaokao; un compañero mayor ya había agotado por completo los 60 puntos de bonificación, y eso le daba mucho miedo. Pero con el apoyo de su familia, eligió la opción de mayor riesgo pero de techo más alto: la rebaja de 60 puntos de Tsinghua.

\begin{warningbox}{Patrón de decisión frente a la incertidumbre}
Esta fue la primera vez que Weng Jiayi enfrentó una decisión importante bajo incertidumbre. Su estrategia fue: en medio del miedo y la incertidumbre, elegir la opción con el \textbf{techo más alto}. Este patrón de decisión volvió a aparecer una y otra vez a lo largo de su carrera profesional: elegir OpenAI en vez de una gran empresa más segura, o dedicarse al código abierto en vez de publicar artículos académicos, sigue siempre la misma lógica.
\end{warningbox}

\subsection{Resumen de la sección}

La trayectoria de crecimiento de Weng Jiayi revela dos rasgos centrales: primero, un \textbf{fortalecimiento autoinducido impulsado por la retroalimentación positiva}, que consiste en obtener retroalimentación positiva de manera constante en las áreas donde destaca, lo que genera un interés propio; segundo, una \textbf{manera de pensar en inversión de largo plazo}, que consiste en sacrificar la eficiencia a corto plazo con tal de acumular ventajas para el futuro. Estos dos rasgos influyeron después de manera profunda en cada una de sus decisiones en Tsinghua, CMU y OpenAI.
\section{Los años en Tsinghua: romper la asimetría de información y el sistema de evaluación}

\subsection{Tareas de acceso abierto: «más útil que donar un edificio»}

Después de ingresar a Tsinghua en 2016, lo primero que hizo Weng Jiayi (翁家翌) que «quedó para la posteridad» fue: \textbf{publicar en acceso abierto en GitHub todas sus tareas y todo el material que había recopilado de generaciones anteriores de estudiantes mayores}.

\begin{importantbox}{La idea de la igualdad de acceso a la información}
Weng Jiayi considera que la asimetría de información es una herramienta de supervivencia muy útil en Tsinghua, pero que todos deberían tener acceso igualitario a esa información. Mucha gente capaz no es hábil para recopilar material; si se les da la oportunidad de tener acceso igualitario a la información, les irá mejor en Tsinghua. Publicar las tareas en acceso abierto no busca que la gente copie, sino que quienes llegan después no tengan que dedicar diez o veinte horas a darle vueltas a un problema sin sentido; ese tiempo se ahorra para hacer algo de más valor.
\end{importantbox}

Este repositorio de GitHub se difundió ampliamente en el Departamento de Ciencias de la Computación de Tsinghua. Weng Jiayi comentó medio en broma: «Detén a cualquier estudiante de generaciones posteriores en Ciencias de la Computación y pregúntale; puede que no conozca a quien donó el edificio de información, pero sí conoce a Weng Jiayi, después de todo, todos sobrevivieron gracias a mis tareas». También se burló de sí mismo: «(mi reconocimiento en Tsinghua) es más útil que donar un edificio».

Esto también generó controversia. Algunos estudiantes de generaciones anteriores se opusieron a publicar las tareas en acceso abierto, pues consideraban que eso haría que los estudiantes posteriores dejaran de pensar de manera independiente. Pero Weng Jiayi insiste en que hizo lo correcto: su objetivo no era que la gente copiara las tareas, sino que las personas capaces pero poco hábiles para recopilar material dejaran de andar exhaustas de un lado a otro y pudieran dedicar más tiempo a lo que realmente querían hacer. De lo contrario, los estudiantes a menudo dedicaban diez o veinte horas a darle vueltas a un problema sin sentido, sin atreverse a preguntarle al ayudante de cátedra, con un beneficio de aprendizaje muy bajo.

\subsection{Iniciación a la investigación: entrar al aprendizaje por refuerzo por casualidad}

En segundo año, Weng Jiayi comenzó a hacer investigación y eligió el laboratorio del profesor Zhu Jun (朱军). En ese momento había tres direcciones posibles: métodos bayesianos, redes generativas antagónicas (GAN) y aprendizaje por refuerzo (RL). Originalmente quería hacer algo relacionado con imágenes usando GAN, pero como no sabía cuál opción correspondía a GAN, eligió por error aprendizaje por refuerzo.

\begin{warningbox}{La aleatoriedad en la elección de investigación}
Weng Jiayi admite que su elección de la dirección de RL fue completamente «random»: no sabía qué número de opción era GAN y eligió al azar. Pero después descubrió que RL era «cosas de videojuegos» y, como le pareció interesante, siguió haciéndolo. Muchas decisiones que parecen cambiar la vida de una persona, en su momento, pudieron haber sido solo un evento aleatorio.
\end{warningbox}

Su primer proyecto de RL fue usar una red neuronal para completar VizDoom, un videojuego de los años noventa, y obtuvo el primer lugar. Pero \textbf{no disfrutó} el proceso: el entorno era demasiado limitado, había que hacer un sobreajuste extremo, la dificultad de ajustar los hiperparámetros era de diez a cien veces mayor que en visión por computadora, y la mayoría de las veces el algoritmo simplemente no funcionaba.

Esta experiencia le dejó dos ideas clave:
\begin{enumerate}
    \item En ese entonces, el cuello de botella central de la investigación en RL no estaba en la innovación algorítmica, sino en el ajuste heurístico de hiperparámetros
    \item Él era más hábil y disfrutaba más construyendo \textbf{la infraestructura que soporta la investigación} que la investigación misma
\end{enumerate}

Weng Jiayi describió lo doloroso que era la investigación en RL en esa época: había que usar métodos muy heurísticos para evitar todo tipo de casos extremos. Cambiar el algoritmo en realidad no era tan fundamental: una tarea que para un humano es simple (por ejemplo, identificar un obstáculo) representaba un nivel de dificultad completamente distinto para la IA. Incluso tuvo una especie de «rechazo fisiológico»: no es que no fuera hábil para ajustar hiperparámetros, sino que \textbf{no le gustaba} hacerlo. Este reconocimiento claro de sí mismo hizo que desplazara su enfoque de «cómo hacer buena investigación» a «cómo hacer que la investigación avance con más fluidez», lo cual dio origen a su camino posterior en infraestructura.

\subsection{Tres direcciones de interés y el despertar del sistema de evaluación}

Además de la IA, durante su paso por Tsinghua, Weng Jiayi también se interesó por la \textbf{computación gráfica} y la \textbf{seguridad de redes}.

En cuanto a computación gráfica, su interés inicial vino de la película de ciencia ficción \textit{Tron: el legado} que vio en secundaria, la cual lo impactó profundamente por sus efectos especiales: «si algún día pudiera crear ese tipo de efectos, o construir mi propio mundo virtual, eso sería la plenitud». En su curso de computación gráfica obtuvo una de las dos únicas calificaciones A+ de toda la clase, inventó un nuevo algoritmo que reducía el número de iteraciones necesarias para la convergencia, y renderizó una imagen con una resolución de 16K sin precedentes en ese momento: antes que él, nadie había renderizado una imagen en 16K. Consideraba que la computación gráfica era «un ejercicio de hacking sobre el mundo real», que le permitía construir, desde su propia perspectiva, las escenas que imaginaba en su mente.

En cuanto a seguridad de redes, le parecía que «ser hacker es sumamente genial»; en su tiempo libre trabajó en muchas cosas relacionadas y corrigió no pocos errores en la red del campus de la escuela. Por ejemplo, él y un estudiante de una generación anterior encontraron una vulnerabilidad que permitía descargar el historial académico por un centavo o incluso gratis (originalmente costaba diez yuanes cada vez); tras descargarlo varias veces, reportaron el error al departamento académico de la escuela.

Pero al final optó por concentrarse en la dirección de IA: «si vas a dedicarte a la investigación, es mejor concentrarse; no se puede tener un pie en cada barco».

\begin{knowledgebox}{El sistema de evaluación de tres indicadores del profesor Zhu Jun}
El asesor de Weng Jiayi, Zhu Jun, propuso un sistema de evaluación distinto del promedio de calificaciones: el valor de un estudiante de Ciencias de la Computación puede medirse con tres indicadores: \textbf{artículos publicados, resultados en competencias y proyectos de GitHub con más de cien estrellas}. Este sistema de evaluación influyó profundamente en Weng Jiayi y lo hizo darse cuenta de que podía crear un valor distinto en la comunidad de código abierto.
\end{knowledgebox}

En cuanto al promedio de calificaciones, la estrategia de Weng Jiayi fue \textbf{invertir el mínimo esfuerzo posible}: calculaba su puntaje actual para asegurarse de alcanzar el estándar que él mismo se había fijado (por ejemplo, un B+ de 87 puntos), y no estaba dispuesto a dedicar tiempo a obtener ni un punto más. Consideraba que el promedio de calificaciones era algo que, tres o cuatro años después, ya no aparecería en el currículum, por lo que no valía la pena invertirle demasiado esfuerzo.

\subsection{Investigación de verano en MILA: el cruce de caminos con Yoshua Bengio}

En el verano de tercer año, Weng Jiayi contactó, a través de su asesor, a Yoshua Bengio (ganador del Premio Turing en 2019) y fue a MILA a hacer investigación de verano. La tarea consistía en construir algo similar a MoE (Mixture of Experts, mezcla de expertos): un router que elegía distintos caminos, aplicado a un modelo de lenguaje.

El resultado de esta experiencia no fue el ideal: en ese momento no había suficiente poder de cómputo ni capacidad de ingeniería para escalarlo, y con unas cuantas tarjetas gráficas simplemente no se podía lograr. Pero, visto en retrospectiva, esta experiencia hizo que Weng Jiayi tuviera contacto simultáneo con RL y NLP (Transformer), lo cual sentó una base doble para su trabajo posterior en RL para LLM en OpenAI.

Weng Jiayi describió la dificultad de ese momento: una persona dedicada a RL lo envió a hacer NLP, y tuvo que dedicar mucho tiempo a introducirse en Transformer y NLP. Lo que logró construir no tuvo buenos resultados; visto en retrospectiva ahora, para que aquello funcionara se necesitaba poder de cómputo y una capacidad de ingeniería muy fuerte para escalarlo, algo que una sola persona con unas cuantas tarjetas gráficas simplemente no podía lograr. «Aunque la dirección fuera correcta, no lo ibas a poder sacar adelante». Esta experiencia de investigación de verano tampoco produjo ningún artículo, lo cual tuvo un impacto considerable en sus solicitudes.

\begin{warningbox}{Los límites de la clarividencia}
Weng Jiayi enfatiza que «opinar después de los hechos no sirve de nada». Aunque, visto en retrospectiva, antes de llegar a OpenAI ya contaba simultáneamente con experiencia en RL y en NLP, en ese momento no podía prever en absoluto el futuro. Para él, en NLP las tareas estaban demasiado dispersas, RL más Transformer se venía abajo, y hacer MoE con unas cuantas tarjetas gráficas no se podía lograr aunque la dirección fuera correcta. \textbf{Muchos reconocimientos que cambian el mundo necesitan suficiente poder de cómputo e ingeniería para poder validarse.}
\end{warningbox}

\subsection{Resumen de la sección}

En sus cuatro años en Tsinghua, Weng Jiayi completó tres transformaciones clave: (1) pasó de una mentalidad de competencia a la idea de acceso abierto e igualdad de información; (2) estableció su propio sistema de evaluación (artículos, competencias, proyectos de código abierto) en lugar del promedio de calificaciones; (3) se dio cuenta de que era más hábil y disfrutaba más haciendo infraestructura que investigación. Estos reconocimientos marcaron el rumbo de su carrera posterior.
\section{El camino del código abierto: Tianshou y tuixue.online}

\subsection{Tianshou (天授): un marco de trabajo de RL hecho a mano en dos semanas}

A principios de 2020, Weng Jiayi (翁家翌) tomó una decisión que cambiaría el rumbo de su carrera: escribir desde cero un marco de trabajo de aprendizaje por refuerzo.

El motivo era sencillo: había escrito mucho código de experimentos de RL y quería integrarlo para que corriera mejor. Primero revisó RLlib, dentro de Ray, y tras un mes descubrió que era demasiado complejo: cientos de miles de líneas de código, demasiadas abstracciones, sin ninguna idea de por dónde modificarlo. Así que decidió \textbf{tirarlo todo y empezar de cero, construyendo a mano un marco de trabajo completamente nuevo}.

\begin{importantbox}{La filosofía de diseño de Tianshou: la consistencia (Consistency)}
Weng Jiayi considera que la característica más importante de un buen proyecto es la \textbf{consistencia}. Los proyectos en los que colaboran muchas personas tienden a degradarse, porque nadie sabe bien qué escribió el otro; los supuestos no se transmiten a tiempo, lo que provoca que el código se copie y pegue y se infle. Tianshou tuvo éxito porque una sola persona lo diseñó e implementó de principio a fin, lo que garantizó la consistencia del conjunto.
\end{importantbox}

La primera versión de Tianshou se completó en solo \textbf{dos semanas}. Weng Jiayi explica que, si las abstracciones están bien planteadas, implementar un algoritmo puede tomar menos de veinte líneas de código. Frente a las cientos de miles de líneas de RLlib, la clave de Tianshou está en un diseño de alto nivel reducido al mínimo: si el usuario quiere modificar alguna función, el diseño ya especifica de antemano «solo se puede modificar aquí».

Lo que Tianshou hizo bien en el fondo fue \textbf{captar la necesidad de los usuarios}: en ese momento, los investigadores de RL necesitaban un marco de trabajo fácil de usar, fácil de modificar y con poco código, que pudieran tomar directamente y, con solo estudiarlo un poco, supieran dónde modificar. Tianshou satisfizo exactamente esa necesidad.

\subsubsection{El desarrollo posterior de Tianshou y su degradación}

Tianshou pasó después a ser un proyecto mantenido por la comunidad de código abierto. Cuando Weng Jiayi entró a trabajar en OpenAI, ya no tuvo tiempo de seguir manteniéndolo, así que transfirió el mantenimiento a la comunidad. La regla que estableció fue: mientras haya una sola persona que tome las decisiones finales, se puede mantener la consistency.

Cinco años después, al mirar atrás, Weng Jiayi reconoce que Tianshou se «degradó un poco», porque su context y el context de quien lo sucedió no coincidían del todo, y el sucesor reescribió parte del código, lo que hizo que el conjunto fuera menos consistent. Pero considera que, a largo plazo, esto es aceptable.

Esta experiencia también lo llevó a tomar, en agosto de 2022, una decisión importante: dejar de desarrollar Tianshou de manera gradual. La razón fue que, tras entrar a OpenAI, se dio cuenta de que existía una brecha enorme entre los toy benchmarks a los que apuntaba Tianshou —como Atari y MuJoCo— y el RL que la industria realmente necesitaba (por ejemplo, la alineación de LLM). «Debería invertir más tiempo en cosas más significativas».

\begin{warningbox}{Sobre la decisión de «no publicar papers»}
Weng Jiayi afirma con claridad que no hizo Tianshou para publicar artículos: «Creo que publicar papers no tiene ningún sentido». En ese momento ya tenía artículos publicados y resultados en competencias, suficientes para sus solicitudes. Lo que quería era un proyecto de código abierto real, con más de tres dígitos de stars, un proyecto de GitHub «como es debido», según el sistema de evaluación de su asesor.
\end{warningbox}

\subsection{tuixue.online: un sistema de consulta de visas durante la pandemia}

Durante la pandemia de 2020, los consulados de Estados Unidos cerraron y los tiempos de espera de visa eran sumamente inciertos. El propio Weng Jiayi también estaba esperando su visa para ir a CMU, así que escribió un rastreador de tiempos de visa y lo publicó como código abierto en forma de un sistema de consulta gratuito: tuixue.online.

La primera versión de este proyecto ni siquiera necesitaba mucha tecnología: se actualizaba a mano una vez de día y una vez de noche. Pero incluso así, tan rudimentario, la demanda ya era enorme. Después de automatizarlo, el número acumulado de clics superó el millón (y más adelante quizá llegó a los diez millones). Cuando terminó la pandemia, los consulados actualizaron sus sitios web, el rastreador original dejó de funcionar y Weng Jiayi tampoco tuvo tiempo de reescribirlo, así que el proyecto cumplió su cometido histórico.

La conclusión de Weng Jiayi sobre este proyecto es: «La tecnología no importa; lo que importa es captar la necesidad». Ni siquiera se necesita tecnología avanzada: incluso una actualización manual, si satisface el punto de dolor real de los usuarios, tiene un valor enorme.

\subsubsection{El código abierto como una forma de caridad}

Weng Jiayi define reiteradamente sus proyectos de código abierto como «caridad». El presentador señala que, en realidad, ni Tianshou ni tuixue.online son proyectos utilitarios: cuando hizo Tianshou, sus solicitudes ya habían terminado, y tuixue es completamente gratuito. Weng Jiayi reconoce que tiene un «impulso interno muy fuerte»: el de crear algo que él considera útil y compartirlo con todos.

Compartió el origen de este impulso: en el último año de preparatoria le surgió de pronto una idea: «Si la vida es un juego, entonces la puntuación final del juego es el número de personas que recuerdan tu nombre en el instante en que mueres». Esta «teoría de la vida como juego» lo ha impulsado desde entonces a buscar impact, y sigue haciéndolo hasta hoy.

El presentador insistió: «¿No sientes que tú también terminas empujado por este criterio?». La respuesta de Weng Jiayi resulta interesante: «Por ahora no. Si me doy cuenta de que pasa eso, puedo cambiarlo; puedo cambiar mi propio criterio de evaluación». No es «esclavo» de ese criterio: aunque durante mucho tiempo en OpenAI no haya tenido ningún proyecto nuevo de código abierto, tampoco le preocupa, porque «el OpenAI model ya es el mejor (proyecto de código abierto)».

\begin{importantbox}{El código como caridad}
Weng Jiayi define tanto Tianshou como tuixue.online como \textbf{proyectos de caridad} (non-profit). Dice: «Hacer proyectos de caridad me da una satisfacción enorme. Comparado con el dinero, el impact me satisface más». Esta idea de «el código como caridad» viene de una «regla de puntuación del juego de la vida» de la que se dio cuenta de pronto en la preparatoria: \textbf{cuantas más personas recuerden tu nombre en el instante en que mueres, más alta es la puntuación}.
\end{importantbox}

\subsection{Resumen de la sección}

Tianshou y tuixue.online comparten el mismo patrón de creación: \textbf{tener una necesidad propia → no existir en el mercado una herramienta buena → construir una a mano → publicarla como código abierto y gratuita para todos}. Ninguno de los dos proyectos es utilitario: Tianshou no se hizo para publicar artículos, tuixue no se hizo para ganar dinero. Pero fueron precisamente estos proyectos «no utilitarios» los que le trajeron a Weng Jiayi el mayor impact y las mayores oportunidades profesionales. La tecnología no importa; lo que importa es \textbf{captar la necesidad}.

\section{De CMU a OpenAI: la elección de carrera}

\subsection{Solicitud de posgrado: ni una carta de recomendación del Premio Turing garantiza un PhD}

La solicitud de Weng Jiayi (翁家翌) para un PhD no salió bien. Aunque contaba con una carta de recomendación de Yoshua Bengio (ganador del Premio Turing), su investigación de verano no produjo ningún artículo, y al final solo obtuvo el Master de CMU en lugar del PhD.

\begin{knowledgebox}{La atmósfera de evaluación de Tsinghua}
En la atmósfera de Tsinghua de aquel entonces, el PhD se consideraba superior al Master. Weng Jiayi admite con franqueza que «sí se sintió un poco decepcionado» y que le tomó un tiempo ajustarse. Pero después llegó a pensar que \textbf{el nivel del título académico no determina el desarrollo a largo plazo}; lo que de verdad importa es si tu experiencia puede corresponder a lo que exige quien contrata. «Si quieres entrar a la industria, entonces hacer un PhD es un desperdicio de vida.»
\end{knowledgebox}

Durante su tiempo en CMU tomó clases en línea desde casa durante un año (debido al COVID), pero en cambio aprovechó ese periodo para desarrollar Tianshou (天授) y tuixue.online, dedicando su energía a cosas de mayor valor a largo plazo.

La experiencia de solicitud de Weng Jiayi ofrece un caso contraintuitivo: una carta de recomendación de un ganador del Premio Turing, más una licenciatura en la carrera de Ciencias de la Computación de Tsinghua, más antecedentes en competencias de OI, y aun así solo se puede obtener el Master. Recuerda que en aquel entonces vio en Zhihu (知乎) a mucha gente discutiendo los resultados de sus solicitudes; algunos, al ver que él tenía el antecedente de una «fuerte recomendación del Premio Turing» pero al final no había obtenido el PhD, se sorprendieron muchísimo: «la competencia es tan intensa».

Pero después Weng Jiayi lo aceptó con calma. Piensa que la distinción de jerarquía entre PhD y Master es una ilusión creada por el ambiente del entorno: «lo que de verdad importa es qué es lo que realmente haces». Empezó a intentar liberarse de este sistema de evaluación; aunque durante la licenciatura todavía no logró liberarse por completo, ya se dio cuenta de que «uno debería crear su propio sistema de evaluación, en lugar de usar el que proporcionan los demás».

\subsection{El Master en CMU: estudios remotos durante la pandemia}

En el otoño de 2020, Weng Jiayi ingresó a CMU, pero debido a la pandemia de COVID y a problemas de visa, pasó todo el primer año tomando clases en línea desde China. Este periodo que parecía «un desperdicio» lo aprovechó al máximo: se concentró en el desarrollo de Tianshou, en el mantenimiento de tuixue.online y en pensar en su futura dirección profesional.

Menciona un cambio de mentalidad importante: frente a narrativas grandiosas como el COVID y la turbulencia política internacional, eligió \textbf{concentrarse en lo que tenía entre manos}: «no hay que estar pendiente todos los días de grandes narrativas internacionales, sino concentrarse en lo que se tiene entre manos; así uno puede sentirse un poco más en paz por dentro».

\begin{knowledgebox}{La decisión de irse al extranjero durante el año de la pandemia}
La generación de Weng Jiayi (que solicitó en diciembre de 2019) coincidió justo con el COVID. En ese entonces, en Tsinghua solo alrededor del 5\% de las personas eligieron ir al extranjero (una caída considerable respecto al 20\% de años anteriores), y enfrentaban múltiples riesgos, como el cierre de visas y la incertidumbre de la situación internacional. Weng Jiayi, en medio de esta incertidumbre, siguió decidido a irse al extranjero, un patrón de decisión idéntico al de cuando en la preparatoria eligió Tsinghua con 60 puntos menos.
\end{knowledgebox}

\subsection{Búsqueda de empleo: DeepSeek vs OpenAI}

Al graduarse de CMU, Weng Jiayi se burla de sí mismo diciendo que «al principio fue bastante descuidado»: envió solicitudes a 18 empresas y al final solo recibió ofertas de Google y de autoML (la empresa del profesor Chen Tianqi (陈天琦)). No quería ir a Google: «ser un tornillo más en una gran empresa, haciendo cosas que no me gustaban tanto, como frontend y backend». Después se lo tomó en serio y consiguió más ofertas.

Las opciones centrales que enfrentaba eran varias empresas:

\begin{table}[H]
\centering
\begin{tabular}{lll}
\toprule
\textbf{Empresa} & \textbf{Área} & \textbf{Resultado} \\
\midrule
OpenAI & RL (equipo de John Schulman) & Aceptó \\
Huanfang (幻方, luego DeepSeek) & RL Infra & Obtuvo oferta \\
NVIDIA & RL & Obtuvo oferta \\
Google & SWE general & Obtuvo oferta \\
FAIR (Meta) & RL & Rechazado por razones de proceso \\
TikTok & -- & Obtuvo oferta \\
\bottomrule
\end{tabular}
\caption{Las opciones de empleo de Weng Jiayi}
\end{table}

\begin{importantbox}{La lógica detrás de elegir OpenAI}
Esta era la época \textbf{before ChatGPT}. Las razones por las que Weng Jiayi eligió OpenAI fueron: (1) OpenAI y DeepMind eran entonces los dos mejores research labs en el campo del RL; (2) quería experimentar en persona cómo se hacía la investigación más avanzada del mundo, en lugar de que unos cuantos PhD la armaran a mano en un «taller artesanal» de la universidad; (3) quería aprender la manera en que la industria hace investigación con metodología.
\end{importantbox}

Sobre Huanfang (antecesor de DeepSeek), Weng Jiayi revela que en ese entonces Huanfang decía que iba a crear un AI Lab (que después se convirtió en DeepSeek). Si no hubiera tenido la oferta de OpenAI, habría elegido Huanfang. Dicho de otro modo, visto en retrospectiva, la elección que enfrentaba en ese momento era, en esencia, \textbf{DeepSeek vs OpenAI}.

\subsection{La entrevista de John Schulman}

Weng Jiayi fue entrevistado y contratado personalmente por John Schulman para OpenAI. La razón central por la que Schulman lo reconoció fue que \textbf{su GitHub era muy hermoso}: una persona con buena capacidad de ingeniería es útil para cualquier proyecto. La última ronda de la entrevista fue un problema abierto muy end-to-end; le dieron tres horas, pero Weng Jiayi lo terminó en dos.

Vale la pena mencionar que esta pregunta de entrevista solo se había puesto a prueba con dos personas: Weng Jiayi y el colega que después trabajó en Codex (Andrew Y.); ambos la aprobaron («la tasa de aprobación fue del cien por ciento»). Weng Jiayi siente una profunda gratitud hacia Schulman, e incluso el día en que Schulman dejó la empresa estuvo triste toda una tarde, apagó la computadora y no hizo nada.

La razón central por la que Schulman apreció a Weng Jiayi merece una reflexión profunda: no fueron los artículos, ni el título académico, sino que \textbf{su GitHub era muy hermoso}. Esto confirma justamente el sistema de evaluación de «estrellas de GitHub de tres dígitos» que propuso el profesor Zhu Jun (朱军): en la industria, el producto de ingeniería visible en verdad resulta más convincente que los artículos académicos.

\begin{knowledgebox}{La opinión de Weng Jiayi sobre hacer un PhD}
Cuando le preguntaron si, al buscar empleo, había considerado hacer un PhD, la respuesta de Weng Jiayi fue tajante: no. «Porque cuando conoces a gente de la industria te das cuenta de que si quieres entrar a la industria, entonces hacer un PhD es un desperdicio de tiempo. Perfectamente puedes usar el Master como trampolín para reunir el nivel que exige la industria a un PhD.» Su estrategia consistió en tener clara la diferenciación: hacer proyectos que le permitieran a la contraparte «elegirte a ti, con Master, en lugar de a otro PhD».
\end{knowledgebox}

\subsection{Resumen de la sección}

La experiencia de búsqueda de empleo de Weng Jiayi encarna varios principios: (1) \textbf{competencia por diferenciación}: usar proyectos de código abierto y capacidad de ingeniería para compensar la carencia de un PhD; (2) \textbf{optimización orientada a la necesidad}: entender qué tipo de persona necesita un AI Lab y luego ajustarse a eso; (3) \textbf{no quedar atado a un sistema de evaluación preestablecido}: ¿el Master es inferior al PhD? Entonces demuéstrate en otras dimensiones.

\section{El trabajo central de Weng Jiayi en OpenAI: la infraestructura de RL para post-entrenamiento}

\subsection{La llegada a OpenAI: un laboratorio a gran escala}

Weng Jiayi (翁家翌) fue el empleado número 280 de OpenAI (ahora la empresa supera los 3,000 empleados). Su primera impresión de OpenAI fue la de «un laboratorio a gran escala»: no tenía tanta metodología como se imaginaba, pero había muchas personas con una intuición de investigación muy fuerte que podían señalar el rumbo.

\begin{importantbox}{La introducción de la «productividad avanzada» de OpenAI}
Un punto de inflexión clave fue cuando Barret, Luke y Liam llegaron los tres desde Google para unirse al equipo de RL de John Schulman. Ellos introdujeron la productividad avanzada de Google, una filosofía: \textbf{no pienses en la idea genial ni en el algoritmo genial; primero deja bien construida la infraestructura, para que la velocidad de iteración pase de treinta veces por semana a trescientas veces por semana}. Weng Jiayi mostró una gráfica: el número de iteraciones por unidad de tiempo es proporcional a la tasa de éxito, lo cual es, en esencia, la manera de proceder del RL (ensayo y error).
\end{importantbox}

\subsection{El nacimiento de ChatGPT: un crecimiento exponencial inesperado}

Weng Jiayi reveló la perspectiva interna del lanzamiento de ChatGPT: la intención original de lanzarlo era solo recopilar algunos datos reales de usuarios; se esperaba que lo usaran quizás diez o veinte mil personas, y que si nadie lo usaba, se cerraría a los cinco días. El resultado fue que la curva de crecimiento de usuarios fue \textbf{exponencial}. La explosión de ChatGPT no fue en absoluto planeada, sino «una serie de reacciones químicas mitad casuales, mitad inevitables».

\begin{knowledgebox}{Las expectativas internas del lanzamiento de ChatGPT}
Las expectativas internas de OpenAI sobre ChatGPT eran muy conservadoras: tras el lanzamiento, quizás lo usarían de diez a veinte mil personas, y si el número de usuarios bajaba después de cinco días, se cerraría. Nadie previó que ocurriría un crecimiento exponencial. Esta historia demuestra que incluso los equipos de vanguardia en IA no pueden predecir con exactitud la reacción del mercado ante un producto. Los productos verdaderamente grandes casi nunca se planean, sino que emergen en la confluencia de lo casual y lo inevitable.
\end{knowledgebox}

\subsection{La construcción de la infraestructura de RL para post-entrenamiento}

El trabajo central de Weng Jiayi en OpenAI fue construir toda la infraestructura de RL de post-entrenamiento. Esta infraestructura se ubica en la capa más alta de toda la pila tecnológica: la capa más orientada al «cliente» (los researchers internos).

Él mismo se fijó una reward function clara: \textbf{maximizar el número de veces que su nombre aparece en el blog de OpenAI}. Para lograr este objetivo, construir infraestructura es más eficaz que hacer un solo proyecto de research, porque la infraestructura se puede escalar: todos usan tu infraestructura, y cada vez que se publica un modelo grande tu nombre tiene que aparecer ahí. El presentador comentó: «De verdad sabes escribirte muy bien tu propia reward».

Weng Jiayi explicó la lógica detrás de esta reward function: si haces un solo proyecto de research, cada vez solo puedes influir en un lanzamiento; pero si haces infraestructura, todos los modelos entrenados con ese conjunto de infraestructura llevarán tu nombre en la lista de colaboradores. Y como se le da bien escribir infraestructura de RL, esta era una «oportunidad muy, muy adecuada».

\begin{importantbox}{El valor central de la infraestructura: corregir errores equivale a la calidad del modelo}
Weng Jiayi planteó un punto de vista incisivo: \textbf{la infraestructura de cada empresa tiene errores en distinto grado, y quien corrija más errores tendrá mejor entrenamiento de modelos.} Incluso especuló que la razón por la que LLaMA no supera a GPT podría ser que la infraestructura de LLaMA tiene más errores. Desde su perspectiva, la esencia de la competencia de vanguardia en IA no es la innovación algorítmica, sino \textbf{la corrección de la infraestructura y la velocidad de iteración}.
\end{importantbox}

\subsection{La siguiente generación de infraestructura: empezar de cero}

Weng Jiayi reveló que OpenAI está \textbf{echando todo abajo para empezar de cero}, reconstruyendo la infraestructura interna de RL. La razón es que la generación anterior de infraestructura ya llevaba más de tres años en operación y había acumulado una gran cantidad de deuda técnica (technical debt). El objetivo de la nueva infraestructura es limpiar la deuda histórica y dar a los researchers una mejor velocidad de iteración.

Reveló que él mismo «en realidad ya no está en la posición más central», pero considera que debería hacer algunas cosas más importantes, y reconstruir la infraestructura es una de ellas. Esto se corresponde con la experiencia de Tianshou: cuando un sistema acumula demasiadas inconsistencias y deuda técnica, la mejor solución puede ser echarlo todo abajo y rediseñar la arquitectura de nivel superior a partir de un nuevo entendimiento.

\begin{warningbox}{El ciclo de vida de la infraestructura y la deuda técnica}
La experiencia de Weng Jiayi muestra que incluso la infraestructura de IA de primer nivel mundial tiene un «ciclo de vida efectivo» de unos tres años. Después, la deuda técnica se acumula hasta un punto en que el costo de aplicar parches supera el costo de reescribir, y entonces es necesario echarlo todo abajo para empezar de cero. Esto es similar al caso de Tianshou: después de que Weng Jiayi se fue, los nuevos mantenedores, al tener un contexto distinto, inevitablemente introdujeron inconsistencias. \textbf{Todo sistema mantenido por personas se degrada poco a poco; lo importante es cuándo y cómo reiniciarlo.}
\end{warningbox}

En el sistema de infraestructura de OpenAI, los researchers no participan en la construcción de la infraestructura: ellos plantean las necesidades, y el equipo de infraestructura las implementa. Al final, el researcher «tal vez solo tenga que cambiar una flag y ya».

\begin{knowledgebox}{La división del trabajo entre researcher y engineer}
Dentro de OpenAI, la infraestructura y la investigación tienen una división de trabajo clara. El researcher se encarga de proponer ideas y necesidades; el engineer se encarga de construir la infraestructura correcta y de iterar rápido. Weng Jiayi considera que «la idea es barata»: las ideas surgen con solo discutirlas con alguien; lo difícil es cuántas ideas válidas se pueden verificar correctamente por unidad de tiempo.
\end{knowledgebox}

\subsection{Resumen de la sección}

La posición de Weng Jiayi en OpenAI es la de «quien vende las palas»: no extrae el oro, pero le da a cada quien que busca oro la mejor pala. La infraestructura de RL de post-entrenamiento se ubica en la capa más alta de la pila tecnológica, con un nicho ecológico sumamente alto. A través de esta posición logró escalar su impacto: cada modelo entrenado con este conjunto de infraestructura lleva su nombre cuando se publica.
\section{Filosofía de la ingeniería y metodología}

\subsection{Enseñarle Engineering a un Researcher es más difícil que al revés}

Weng Jiayi (翁家翌) cita a un colega (un PhD que trabajó en un RL framework reconocido):

\begin{importantbox}{La asimetría entre Engineering y Research}
«Enseñarle a un researcher a hacer bien engineering es mucho más difícil que enseñarle a un engineer a hacer bien research.» En la era actual de la IA, el valor de la capacidad de ingeniería supera al de la capacidad académica, porque la intuición de research puede acumularse trabajando mucho tiempo en la industria, mientras que una capacidad de ingeniería sólida requiere un entrenamiento deliberado y prolongado.
\end{importantbox}

\subsection{La velocidad de iteración lo determina todo}

La idea central que Weng Jiayi (翁家翌) enfatiza una y otra vez es que \textbf{el número de iteraciones correctas por unidad de tiempo} es el factor determinante del éxito en la investigación de IA. En concreto:

\begin{enumerate}
    \item Corrección de la infra: una infra sin bugs significa que el resultado de cada experimento es confiable
    \item Velocidad de iteración: pasar de 30 iteraciones por semana a 300 iteraciones por semana hace que la tasa de éxito crezca de forma lineal
    \item Rendimiento en la corrección de bugs: cuántos bugs se pueden corregir por unidad de tiempo, cuántas iteraciones correctas se pueden lograr
\end{enumerate}

\begin{warningbox}{La innovación algorítmica no es el cuello de botella}
Weng Jiayi (翁家翌) considera que el cuello de botella de la frontera actual de la IA \textbf{no está en la innovación algorítmica}. «Si corriges todos los bugs, es posible que el algoritmo funcione bien sin necesidad de cambiar nada.» Muchos problemas que parecen requerir un nuevo algoritmo para resolverse en realidad pueden deberse solo a bugs de la infra. Esto contrasta claramente con el paradigma académico, centrado en proponer nuevos algoritmos.
\end{warningbox}

\subsection{Consistencia del código y deterioro de los proyectos}

Weng Jiayi (翁家翌) considera la \textbf{consistencia (consistency)} como el primer principio tanto de los proyectos de software como de la gestión de organizaciones. La raíz del deterioro de los proyectos es la falta de consistencia: cuando varias personas colaboran, el context de cada una es distinto, los supuestos no pueden transmitirse, y esto provoca que el código se infle y su calidad se deteriore.

Él considera que dirigir una empresa y mantener un código base son esencialmente lo mismo: ambos requieren mantener la consistency. Si no hay consistencia, el código base se vuelve inflado y la estructura organizacional también se vuelve inflada. La solución es que \textbf{la información fluya sin obstáculos}: que las decisiones de arriba lleguen abajo sin pérdida, y que el progreso de abajo llegue arriba sin pérdida.

Weng Jiayi (翁家翌) lleva este planteamiento hasta el extremo: en teoría, el \textbf{context sharing debería ser reemplazado por un agent con un context de longitud infinita}. Esto porque el context del cerebro humano es limitado y no puede almacenar al mismo tiempo toda la información de una organización entera. Pero la IA sí puede. En el futuro, tal vez cada empresa tenga un agent de context infinito de este tipo que funja como CEO, encargado de todo el sharing y las decision: «Probablemente no haya un decision maker más adecuado que un agent así.» Esta perspectiva unifica su filosofía de la ingeniería, sus convicciones sobre la IA y sus reflexiones sobre la gestión organizacional.

\subsection{Resumen de la sección}

La filosofía de la ingeniería de Weng Jiayi (翁家翌) puede resumirse en tres principios centrales: (1) \textbf{prioridad a la consistencia}: que una sola persona se encargue de todo, lo cual, aunque no favorece la escalabilidad a largo plazo, garantiza la consistencia; (2) \textbf{prioridad a la velocidad de iteración}: no se busca un algoritmo genial de una sola vez, sino aproximarse a la solución óptima mediante prueba y error rápidos; (3) \textbf{prioridad a la corrección}: corregir bugs es más importante que escribir un nuevo feature.


\section{Perspectivas de la industria: academia vs. industria}

\subsection{¿Todavía vale la pena cursar un PhD?}

Weng Jiayi (翁家翌) dio un consejo claro a los jóvenes que enfrentan esta disyuntiva en 2025:

\begin{importantbox}{El consejo profesional de Weng Jiayi}
\textbf{Si quieres entrar a la industria, entonces cursar un PhD es un desperdicio de vida.} Puedes usar perfectamente una maestría o incluso la licenciatura como trampolín, y competir en igualdad de condiciones con los candidatos a PhD acumulando experiencia diferenciada en proyectos (sobre todo en infra). «El propósito principal de contratar es contratar a gente que sirva, que pueda hacer el trabajo.»
\end{importantbox}

Él considera que la academia se está reestructurando. Lo que más necesita un AI Lab en este momento es talento de infra: infra es un «pozo sin fondo», la gente con intuición de research ya es de por sí muy poca (los que llevan más de tres años en la industria se cuentan con los dedos de una mano), y todo el cuello de botella restante está en infra.

Weng Jiayi sugiere que el razonamiento de los jóvenes debería ser: primero averiguar con claridad qué tipo de persona necesita realmente un AI Lab, y después ajustarse a eso. Si lo que más necesitan es gente de infra, hay que hacer más trabajo de infra, incluso sin un PhD degree no hay problema: «lo más importante es ver si tu experiencia hace match, si sirve de algo». Él enfatiza que si eres un new grad (recién egresado) con una experiencia muy alineada con el trabajo objetivo, «eso puede equivaler a varios años de experiencia laboral».

\begin{knowledgebox}{La evolución del sistema de evaluación}
La experiencia de Weng Jiayi muestra la evolución del sistema de evaluación de la escuela a la industria: Tsinghua usa el GPA → el tutor usa artículos/concursos/estrellas de GitHub → OpenAI usa las contribuciones a infra y el número de lanzamientos de modelos. Los criterios de evaluación son distintos en cada etapa; lo importante es \textbf{identificar con anticipación los criterios de la siguiente etapa}, en lugar de sobreoptimizar en la etapa actual. Esto en sí mismo es una estrategia de «invertir en el futuro».
\end{knowledgebox}

\subsection{La desconexión entre la investigación de RL en la academia y la industria}

Weng Jiayi señala la enorme brecha entre la academia y la industria en el campo de RL:

\begin{table}[H]
\centering
\begin{tabular}{ll}
\toprule
\textbf{Academia} & \textbf{Industria} \\
\midrule
Hacer overfit en unas cuantas tasks de Atari / MuJoCo & Usar RL para resolver problemas reales (como la alineación de LLM) \\
Ver quién obtiene la puntuación más alta en el step 100K & Ver qué modelo funciona mejor en escenarios reales de usuarios \\
Buscar publicar artículos con nuevos algoritmos & Buscar la corrección de infra y la velocidad de iteración \\
Experimentos pequeños con unas cuantas GPUs & Entrenamiento distribuido a gran escala \\
\bottomrule
\end{tabular}
\caption{Comparación entre academia e industria en la investigación de RL}
\end{table}

Después de darse cuenta de esto en agosto de 2022, Weng Jiayi dejó de desarrollar poco a poco Tianshou (天授), porque Tianshou seguía enfocado en toy benchmarks, y él debía dedicarse a algo con más sentido. Dice que en ese momento tuvo claro que «debería dedicar más tiempo a algo con más sentido»: en cuanto vio con claridad la diferencia esencial entre el RL académico y el RL industrial, seguir manteniendo un marco de trabajo enfocado en toy benchmarks perdió el atractivo.

Weng Jiayi también citó una frase de un colega suyo para resumir esta brecha:

\begin{importantbox}{Idea is cheap --- lo que cuesta caro es la capacidad de validar ideas}
«Las ideas son baratísimas. Lo que hay que hacer es lograr validar la mayor cantidad posible de ideas válidas por unidad de tiempo, y que sea con la infra correcta, el resultado correcto y la iteración rápida.» Quien de verdad piensa puede ser alguien como Alec Radford, que ha estado trabajando en esto desde GPT-1; su intuición de research es más útil que el razonamiento de un PhD común, así que para las ideas basta con discutirlas con él. Todo el cuello de botella restante está en execution.
\end{importantbox}

\subsection{Resumen de la sección}

Hoy, con el rápido desarrollo de la IA, la ruta académica tradicional (PhD → Paper → puesto docente) se está reevaluando. La experiencia de Weng Jiayi muestra que la capacidad de ingeniería, la experiencia diferenciada en proyectos y una comprensión precisa de las necesidades pueden determinar el desarrollo profesional de una persona más que el nivel de estudios y el número de artículos publicados.

\section{La cultura organizacional de OpenAI y el panorama competitivo}

\subsection{Densidad de talento y estructura organizacional}

Weng Jiayi (翁家翌) está de acuerdo con lo que dijo Sam Altman: «En un equipo pequeño con una densidad de talento extremadamente alta, no se tolera ningún desempeño mediocre.» Una densidad de talento alta permite que surjan de manera espontánea cosas inesperadas.

\begin{knowledgebox}{El mecanismo de circulación de información en OpenAI}
La clave para que OpenAI mantenga su capacidad de innovación es que \textbf{la información circule sin obstáculos}. Sam Altman tiene un asistente de investigación dedicado que lo ayuda a mantenerse al tanto de los avances internos más recientes; Greg Brockman ha participado en casi toda la infra subyacente. El conocimiento profundo que el liderazgo tiene de los detalles técnicos garantiza que la información sea consistente entre el nivel de decisión y el nivel de ejecución.
\end{knowledgebox}

Al pasar de 280 personas a más de 3,000, el reto central que enfrenta OpenAI es cómo conservar la eficiencia de innovación de un equipo pequeño. Weng Jiayi considera que la probabilidad está bajando, pero no demasiado: lo importante es poder separar equipos pequeños dedicados a la investigación, además de simplificar la estructura organizacional y eliminar los meeting injustificados.

\subsection{La perspectiva interna sobre el despido de Sam Altman}

Sobre el episodio de noviembre de 2023 en que la junta directiva destituyó a Sam Altman, Weng Jiayi ofrece una perspectiva interna: las personas que hacían el trabajo en la base \textbf{no tenían ninguna idea de lo que estaba pasando}, fue un surprise total. La junta directiva carecía de transparencia hacia los empleados y el proceso de decisión no fue transparente. La causa central fue la \textbf{desconfianza} que Ilya y otros miembros de la junta tenían hacia Sam: no desconfianza de alguna decisión técnica en particular, sino desconfianza de la persona misma.

La exigencia central de Weng Jiayi es la \textbf{estabilidad organizacional}: «que no vuelva a pasar que la empresa esté a punto de quebrar». La estabilidad de la estructura organizacional favorece avanzar con rapidez. Desde su punto de vista, tanto la mayor oportunidad como el mayor reto para lograr la AGI se resumen en una sola palabra: \textbf{ejecución}: «ejecutar en la dirección correcta; basta con poder ejecutar.»

\subsection{La competencia de DeepSeek y los retos de OpenAI}

Weng Jiayi analiza con franqueza las ventajas competitivas de DeepSeek:

\begin{warningbox}{La desventaja de escala de las empresas grandes}
Weng Jiayi considera que la eficiencia de DeepSeek es mucho mayor que la de OpenAI: su base de código es pequeña, el costo de comunicación es bajo y se concentra en un use case específico. OpenAI, en cambio, tiene que tomar en cuenta muchos use cases a la vez y hacer trade off en cada aspecto; al crecer la organización, enfrenta inevitablemente una caída de eficiencia. Pero OpenAI tiene ventaja en otras dimensiones, como la retroalimentación de los usuarios; eso también es un trade off. \textbf{Todas las empresas se vuelven más lentas: lo que importa es quién se vuelve menos lenta.}
\end{warningbox}

Weng Jiayi revela que, después de que DeepSeek liberó su código, dentro de OpenAI sí se reevaluó la estrategia de código abierto. John Schulman incluso le preguntó si convenía liberar el RL Infra como código abierto, pero considerando los intereses de la empresa, Weng Jiayi pensó en ese momento que no era muy adecuado. Sin embargo, dijo que si hubiera recursos ilimitados, «por supuesto que le daría mucho gusto» liberarlo como código abierto.

Aquí entra un problema de teoría de juegos: si OpenAI libera como código abierto su tecnología más avanzada, otras empresas la alcanzarían de inmediato, y entonces OpenAI podría no conseguir financiamiento ni tener quién siga inyectándole recursos. Weng Jiayi considera que «no todos tienen las mismas intenciones»: aunque OpenAI quisiera liberar su tecnología en código abierto para que la AGI beneficie a toda la humanidad, también hay quien solo quiere aprovechar eso para ganar dinero. Para prevenir esa situación, OpenAI se ve obligada a mantener su código cerrado.

\begin{knowledgebox}{El dilema de teoría de juegos del código abierto}
Weng Jiayi señala la contradicción central de la estrategia de código abierto: en teoría existe un camino —liberar el código y recibir retroalimentación de la comunidad permitiría lograr la AGI de mejor manera. Pero llevarlo a la práctica es muy difícil, porque (1) si eres el número uno y liberas tu código, otro se vuelve el número uno de inmediato; (2) a otros les basta con entrenar un poco más para superarte; (3) eso hace que no consigas financiamiento. Este es un dilema del prisionero típico: la racionalidad individual produce un resultado colectivo subóptimo.
\end{knowledgebox}

\subsection{Sobre el «Open» de OpenAI}

Weng Jiayi ofrece una interpretación singular del «Open» de OpenAI:

\begin{importantbox}{El Open de OpenAI: dirigido a la gente común}
El Open de OpenAI no es un Open hacia otras empresas de modelos grandes, sino \textbf{un Open hacia la gente común}: al precio más barato posible (incluso con un tier gratuito), que todos puedan tener acceso a la tecnología más avanzada. «Si le das a alguien los pesos desnudos de un modelo, la gente común tampoco sabe cómo usarlos.» El verdadero Open es hacer que la tecnología esté al alcance de la mano.
\end{importantbox}

\subsection{Resumen de la sección}

La contradicción central que enfrenta OpenAI es la tensión entre \textbf{el crecimiento en escala y la eficiencia de innovación}. La consistencia en la circulación de información —ya sea en la base de código, la estructura organizacional o la comunicación entre personas— es la solución que Weng Jiayi enfatiza una y otra vez. Incluso imaginó un escenario extremo: en el futuro podría haber un AI agent con context ilimitado que funja como CEO, encargado de todo el context sharing y el decision making.

\section{Perspectivas futuras y la frontera de la IA}

\subsection{¿RL for LLM todavía necesita un Breakthrough?}

La respuesta de Weng Jiayi (翁家翌) es que \textbf{es posible}, pero la tarea más urgente es primero exprimir al máximo los métodos y los recursos de cómputo existentes. El estado actual es que «todavía no se ha hecho scale up por completo»: primero hay que hacer hill climb, ver cuál es el límite superior de los métodos existentes, y después considerar si se necesita un nuevo paradigma.

Él enfatizó una idea clave: \textbf{no se puede usar el estado actual para predecir qué sucederá después}. Podría haber un nuevo paradigma de RL, podría haber un nuevo paradigma de post-training, cualquiera de los dos es posible. Cada día se enfrenta a retos desconocidos.

\subsection{¿Dónde está el mayor cuello de botella?}

Weng Jiayi considera que el mayor cuello de botella de los modelos grandes en el futuro sigue siendo la infra:

\begin{enumerate}
    \item El \textbf{rendimiento} para reparar la infra: cuántos bugs se pueden corregir por unidad de tiempo
    \item Cuántas veces se puede \textbf{iterar correctamente} por unidad de tiempo
    \item Estos dos indicadores determinan la eficiencia productiva de todo el equipo y pueden potenciar todo el resto del trabajo
\end{enumerate}

«Ya sea el algoritmo o el entorno: si arreglas todos los bugs, es posible que el algoritmo funcione muy bien sin necesidad de cambiar nada.»

\subsection{Si la IA pudiera resolver un gran problema del mundo}

Cuando le preguntaron qué gran problema del mundo le gustaría que la IA resolviera, la respuesta de Weng Jiayi fue sorprendente: \textbf{cómo predecir el futuro}, no predecir cómo cae una taza, sino predecir toda una vida, la configuración del mundo, todo. Esta respuesta es coherente con su visión determinista del mundo: si el mundo es un proceso de Márkov determinista, en teoría debería poder predecirse.

Pero al mismo tiempo también es consciente de que una herramienta así podría ser un desastre: «Si se consigue una máquina capaz de predecir el futuro, eso sería en realidad un desastre para cada persona: provocaría el colapso de todo el sistema de valores.» Incluso considera que, si llegara a desarrollarse un modelo de IA así, «la mejor opción sería destruirlo, para que nunca llegara a salir».

\subsection{La rutina diaria y el cuidado del cuerpo de Weng Jiayi}

Weng Jiayi reveló cómo es su estado de trabajo en OpenAI. Considera que su trabajo de mantenimiento cotidiano «en realidad no requiere tanto coeficiente intelectual»: basta con hacer lo correcto en la dirección correcta. También compartió un detalle muy contrastante: en Tsinghua (清华) reprobó los 3,000 metros en la clase de educación física, pero ahora tiene el hábito de correr 3,000 metros dos veces por semana: «Antes que nada, hay que asegurarse de que el cuerpo esté sano.» Esto es una extensión de su filosofía de «invertir en el futuro» aplicada al cuidado del cuerpo.

\begin{warningbox}{La humildad y la autoconciencia de un genio}
Weng Jiayi enfatizó varias veces que él no es especial: «Si me sustituyeran por cualquier otra persona, si esa persona tuviera mi context, debería poder desempeñarse perfectamente bien también». Considera que tiene mucha suerte de estar en esta posición, pero que esto «lo podría hacer cualquier ser humano normal». Esta humildad podría ser sincera: él atribuye su éxito más a \textbf{el lugar correcto} (OpenAI + RL Infra) y a \textbf{un context suficiente}, y no al talento personal.
\end{warningbox}

\subsection{Resumen de la sección}

Para Weng Jiayi, el futuro de la IA está lleno tanto de certidumbre (la ruta del scale up es clara, la infra necesita optimizarse constantemente) como de incertidumbre (un nuevo paradigma podría surgir en cualquier momento). La estrategia que elige para enfrentarlo es la de siempre: hacer lo correcto en la dirección correcta, dejar la infra bien construida y permitir que la iteración rápida lo resuelva todo.

\section{Visión del mundo: determinismo, libre albedrío y el sentido de la vida}

\subsection{Un proceso de Márkov determinista}

En la parte filosófica de la entrevista, Weng Jiayi (翁家翌) mostró una faceta muy distinta de la de los temas técnicos. Sostuvo con firmeza:

\begin{importantbox}{La visión del mundo de Weng Jiayi: determinismo}
\textbf{Vivimos dentro de un proceso de Márkov determinista.} Todo puede predecirse: lo que piensas, la siguiente palabra que vas a decir, todo quedó fijado en el instante mismo del Big Bang. El ser humano no tiene libre albedrío, y el mundo macroscópico no tira los dados (la aleatoriedad de la mecánica cuántica a nivel microscópico puede explicarse como una «modificación de las líneas de mundo en segundo plano»).
\end{importantbox}

Él reconoce que se trata de una «visión del mundo bastante pesimista» y que, en el fondo, tampoco quisiera aceptarla, como si él mismo se hubiera convertido en un átomo simulado. Pero, con base en su experiencia personal, considera que en efecto es un hecho. Afirma que «ya lo he verificado innumerables veces», aunque no puede revelar los detalles de esa experiencia personal.

El entrevistador insistió: ¿y la incertidumbre de la mecánica cuántica? La respuesta de Weng Jiayi fue que «a nivel macroscópico no se tiran los dados, a nivel microscópico sí»; pero la aleatoriedad microscópica puede explicarse como una «modificación de algunas líneas de mundo en segundo plano». Reconoció que «se puede pensar que lo que digo son puras tonterías», pero aun así mantuvo su postura.

Esta visión del mundo genera una tensión interesante con su trabajo en RL: el supuesto central de RL es que el agente puede modificar la recompensa futura eligiendo distintas acciones, pero si el mundo es determinista, entonces la «elección» misma del agente también estaría predeterminada. ¿Es Weng Jiayi consciente de esta tensión? Su respuesta lo sugiere: «hay quien quiere desarrollar este tipo de modelos (que predicen el futuro) solo para averiguar cuál es la regularidad que rige este mundo por debajo».

\subsection{La paradoja de invertir en el futuro}

El entrevistador señaló con agudeza una contradicción en la visión del mundo de Weng Jiayi: si todo está determinado, se llegará a ese punto tanto si inviertes en el futuro como si no, entonces ¿para qué invertir?

La respuesta de Weng Jiayi fue: «de todos modos es mejor invertir un poco». Pero al insistírsele en que «bajo un proceso determinista tu inversión no tendría ningún efecto», reconoció: «invertir en el futuro también podría ser algo determinista; el que inviertas o no tampoco es tu libre albedrío».

\subsection{Un periodo de desorientación vital}

Al final de la entrevista, Weng Jiayi compartió con franqueza su estado actual:

\begin{knowledgebox}{El estado actual de Weng Jiayi}
Weng Jiayi atraviesa un periodo de desorientación en cierto punto de su carrera. El trabajo en RL Infra, que alguna vez le gustó mucho, se ha ido volviendo cada vez más determinista con el paso del tiempo. Dice que «alguna vez tuve claro qué quería, pero ahora ya no lo tengo claro». Su meta a corto plazo es retirarse anticipadamente, acumular capital suficiente y luego dedicar tiempo a encontrar lo que de verdad quiere hacer. Respecto a sí mismo dentro de diez años, su única expectativa es «hacer lo que quiera hacer en ese momento, con los recursos y la capacidad suficientes».
\end{knowledgebox}

\subsection{Resumen de la sección}

Las reflexiones filosóficas de Weng Jiayi le dan profundidad a toda la entrevista. Un ingeniero que construyó la infraestructura central de OpenAI resulta, en el terreno filosófico, un determinista absoluto. Considera que la AGI es «un hecho consumado» y que todo lo demás es «algo bastante determinista», que «ya se le ve el final». Pero, paradójicamente, esa certeza es justo lo que lo hace sentirse desorientado: si todo está determinado, ¿dónde queda el sentido del esfuerzo?

Respecto a sí mismo dentro de diez años tiene una única expectativa: «hacer lo que quiera hacer en ese momento, con los recursos y la capacidad suficientes». El entrevistador insistió: «¿no vas a intervenir en lo que él quiera en ese momento?». Weng Jiayi respondió: «porque las ideas cambian; puede que no importe lo que tú quieras». Y añadió: «lo único que puedo hacer ahora es invertir en ese yo de entonces, o dicho de otro modo, invertir en el futuro, para que él tenga la libertad de elegir».

Su mensaje final fue: \textbf{«explorar, al final de cuentas, qué es lo que uno mismo quiere: esta pregunta vale la pena pensarla toda la vida».}

\section{Síntesis y ampliación}

\subsection{Ideario central de Weng Jiayi (翁家翌)}

A lo largo de toda la entrevista a fondo de dos horas, el ideario central que exhibe Weng Jiayi puede resumirse en los siguientes puntos:

\begin{enumerate}
    \item \textbf{Invertir en el futuro}: desde aprender en la secundaria, por adelantado, las matemáticas de la preparatoria, hasta hacer en la universidad proyectos de código abierto sin fines utilitarios, hasta construir en OpenAI infra de propósito general en lugar de un solo research project: el pensamiento de largo plazo lo atraviesa todo.

    \item \textbf{Romper la brecha de información}: código abierto de sus tareas, código abierto de Tianshou (天授), código abierto de tuixue: él siempre ha estado haciendo el trabajo de la equidad informativa. «El código es caridad» no es una consigna, sino una práctica.

    \item \textbf{Prioridad a la consistencia}: ya sea en la base de código, en la estructura organizacional o en el desarrollo personal, consistency es el primer principio. La corrupción proviene de la inconsistencia.

    \item \textbf{Un sistema de evaluación propio}: no dejarse atar por estándares externos como el GPA, el PhD o el número de artículos, sino encontrar las dimensiones de evaluación que le convienen a uno mismo (impact, GitHub star, número de apariciones en el OpenAI Blog).

    \item \textbf{Impulsado por la necesidad}: Tianshou, tuixue, RL Infra: cada proyecto comenzó a partir de una necesidad real. La tecnología no importa; lo que importa es captar la necesidad.
\end{enumerate}

\subsection{Enseñanzas para los jóvenes}

\begin{importantbox}{Resumen de la experiencia de Weng Jiayi}
(1) Primero hay que tener claro qué necesita tu usuario objetivo, y después decidir qué hacer, ya sea para buscar empleo o para emprender.\\
(2) La diferenciación importa más que el nivel académico: un Master puede vencer a un PhD por medio de un excelente proyecto de ingeniería.\\
(3) Idea is cheap, execution is everything: la capacidad de iterar correctamente por unidad de tiempo es la verdadera ventaja competitiva central.\\
(4) Encuentra tu propio sistema de evaluación, pero no te conviertas en su esclavo.\\
(5) Haz proyectos de tipo «caridad»: el impact y la satisfacción que se obtienen en el proceso de ayudar a otros suelen ser más duraderos que la recompensa monetaria.
\end{importantbox}

\subsection{Lecturas adicionales}

\begin{itemize}
    \item Repositorio de GitHub de Tianshou: \url{https://github.com/thu-ml/tianshou}
    \item Repositorio de tareas de código abierto de Weng Jiayi en Tsinghua
    \item El OpenAI Blog oficial: el nombre de Weng Jiayi puede encontrarse en la contributor list de cada model release
    \item El informe técnico de DeepSeek: para conocer el panorama competitivo que menciona Weng Jiayi
    \item Los artículos y las charlas de John Schulman relacionados con RL
\end{itemize}

%% --- Fin del contenido --- %%

\end{document}
